\section{Auxiliary adjoint ``heavy quark" approach}\label{sec:aux}

Let us start with the gluon quasi-PDF defined in Ref.~\cite{Ji:2013dva}
\begin{align}\label{gluonpdfop}
\tilde g(x, P^z)=\int\frac{dz}{2\pi x n\cdot P}e^{i x z n\cdot P}\langle P|O(z,0)|P\rangle,\non\\
%\label{gluonpdfop}
O(z_2,z_1)=F_a^{z\mu}(z_2)L_{ab}(z_2, z_1)F_{b,\mu}^z(z_1).
\end{align}
$F_a^{\mu\nu}=\partial^\mu A_a^{\nu}-\partial^\nu A_a^\mu-g f_{abc}A_b^\mu A_c^\nu$ is the field strength tensor. Unless otherwise stated, the quantities defined in this section are all bare ones. The index $\mu$ can be summed either over transverse components, as in the gluon PDF~\cite{Collins:1984xc}; or over all Lorentz components. Since $F^{zz}=0$ and $F^{zt}$ gets power suppressed in the large momentum limit, both summations can be used to define the gluon quasi-PDF and they fall into a universality class of the gluon quasi-PDF operators~(for discussion on universality class of the gluon polarization operator, see Ref.~\cite{Hatta:2013gta}). $P^\mu=(P^0, 0,0, P^z)$ is the momentum of the external hadron, and $n^\mu=(0,0,0,1)$ is a spacelike vector specifying the direction of the gauge link $L(z, 0) = {\cal P}\exp\left(-ig\int^z_0 d\lambda\,n^\mu  A_\mu(\lambda n)\right)$ with $A_\mu=A_{\mu, a}T_a$ being defined in the adjoint representation.

Alternatively, the gluon quasi-PDF can be defined using fundamental gauge links $U(z_2,z_1)$~\cite{Dorn:1981wa}
\beq\label{fundrep}
\hspace*{-1em}O(z_2,z_1)=2\,{\rm Tr}[F^{z\mu}(z_2)U(z_2,z_1)F_{z\mu}(z_1)U(z_1,z_2)],
\eeq
where $F^{\mu\nu}=F_{a}^{\mu\nu}t_a$ and $t_a$ is the generator in the fundamental representation.
While Eq.~(\ref{gluonpdfop}) facilitates the study of renormalization, Eq.~(\ref{fundrep}) is more straightforward for lattice implementations. We will mainly focus on the definition in Eq.~(\ref{gluonpdfop}) below, but the results also apply to Eq.~({\ref{fundrep}}).

To renormalize $O(z_2,z_1)$, we introduce an auxiliary adjoint ``heavy quark" field, denoted as $\mathcal Q$, into the QCD Lagrangian, following our study of the quark quasi-PDF~\cite{Ji:2017oey}. This ``heavy quark" has trivial spin degrees of freedom. The Lagrangian with the auxiliary field reads
\beq\label{efflag}
\mathcal L=\mathcal L_{\rm{QCD}}+\overline{\mathcal Q}(x)(i n\cdot D-m)\mathcal Q(x),
\eeq
where $D_\mu=\partial_\mu+ig A_{\mu}$ is the covariant derivative in the adjoint representation, and a ``residual mass" is introduced for the auxiliary field $\mathcal Q$, for reasons that will become clear later.

With the auxiliary ``heavy quark" $\mathcal Q$, the Wilson line can be replaced by a product of two such fields~\cite{Ji:2017oey}
\begin{align}\label{Eq:repl}
{\cal O}(z_2, z_1) &= F_a^{z\mu}(z_2){\mathcal Q}_a(z_2) \overline{\mathcal Q}_b(z_1) F^z_{b, \mu}(z_1)\non\\
&= J_1^{z\mu}(z_2) {\overline J}_{1, \mu}^z(z_1).
\end{align}
The renormalization of the non-local operator in Eq.~(\ref{gluonpdfop}) now follows from that of the product of two local composite operators (LCOs) in Eq.~(\ref{Eq:repl}).
